\chapter*{Abstract}
\addcontentsline{toc}{chapter}{Abstract}

California's utilities reported 23,277~MW of data center energization requests to the California Energy Commission (CEC) in December 2025, of which 20,677~MW are located in the state: 5,086~MW with signed agreements, 6,987~MW in active applications and 8,604~MW of inquiries. That is 40 percent of the record CAISO peak and about twenty times the roughly 1,000~MW of data center load that exists. This report asks how much of it is likely to be served by 2030, whether California's supply can keep up, what a flexible load could do about the difference, and how well the state's way of counting compares with the other regimes in use. It is an independent replication and audit of the CEC's data center forecast method, built on public data with every input recorded in a manifest and every figure traceable to its files.

The CEC's published confidence levels and utilization factor, with SVP's requests exempt as the Commission states, reproduce the CEC's own 2040 endpoints exactly and its 2030 components within 2 to 6 percent. Applied symmetrically to the supply side, a completion model estimated on 925 unit-vintages of California's own planned generators since 2016 removes a third of the 20~GW planned list for 2030. A 10,000-draw Monte Carlo over both sides puts the 2030 gap between demand growth and supply growth at 50~TWh (19 to 78) and 7.2~GW at the net peak (4.2 to 10.1), with data centers 18~TWh and 2.3~GW of it; the proposal's upper bound, every requested megawatt operating flat, gives 202 to 242~TWh and 25 to 30~GW. Imports, the state's own load forecast, hydro and Diablo Canyon drive the energy gap, not the data center assumptions. Replicating the Duke Nicholas Institute's method on CAISO hourly load gives 3.8~GW of curtailment-enabled headroom at 0.5 percent curtailment, twice the probability-weighted data center addition and half the weighted peak gap. Under six counting regimes the same requests are 20,677~MW (California and Texas rules), 5,086~MW (PJM and the Texas forecast rule), about 2,400~MW at ERCOT's realized phase-transition rates and about 1,000~MW as measured load. The data are frozen on September 22, 2026, the repository is tagged, and the processed outputs ship with the report.
